\section{First-Fit 空闲链表}

\subsection{\texttt{block\_header\_t} 结构}

First-Fit 后端把堆区域切成若干块，每块前面有小 header，所有块串成双向链表：

\codefile{src/kernel/mm/heap\_first\_fit.c}
\begin{lstlisting}[style=cstyle]
typedef struct block_header {
    size_t size;                   /* 用户数据区大小（不含 header） */
    int is_free;                   /* 0=已分配, 1=空闲 */
    struct block_header* next;
    struct block_header* prev;
} block_header_t;
#define HEADER_SIZE  ((sizeof(block_header_t) + 7u) & ~7u)  /* 16B */
#define ALIGN8(x)   (((x) + 7u) & ~7u)
\end{lstlisting}

\begin{figure}[H]
  \centering
  % figures/ch10/fig03-block-header-layout.tex — auto-extracted from chapter source
\begin{tikzpicture}[
  cell/.style={draw, minimum height=0.7cm, font=\ttfamily\small, anchor=west},
]
  \node[font=\scriptsize, anchor=south] at (0.2, 0.9) {\texttt{block}};
  \draw[-{Stealth[length=4pt]}, thick, blue!60!black] (0.2, 0.9) -- (0.2, 0.7);
  \node[cell, minimum width=1.6cm, fill=red!12]  at (0, 0) {\small size};
  \node[cell, minimum width=1.4cm, fill=red!12]  at (1.6, 0) {\small free};
  \node[cell, minimum width=1.6cm, fill=red!12]  at (3.0, 0) {\small next};
  \node[cell, minimum width=1.6cm, fill=red!12]  at (4.6, 0) {\small prev};
  \node[cell, minimum width=4.2cm, fill=green!12] at (6.2, 0)
    {\small 用户数据区};
  \draw[{|}-{|}, thick] (0, -0.6) -- (6.2, -0.6);
  \node[font=\scriptsize, below] at (3.1, -0.6) {HEADER\_SIZE = 16B};
  \draw[{|}-{|}, thick] (6.2, -0.6) -- (10.4, -0.6);
  \node[font=\scriptsize, below] at (8.3, -0.6) {\texttt{size} 字节};
  \node[font=\scriptsize, anchor=south, text=green!40!black] at (6.4, 0.9)
    {\texttt{kmalloc} 返回此处};
  \draw[-{Stealth[length=4pt]}, thick, green!50!black] (6.4, 0.9) -- (6.4, 0.7);
\end{tikzpicture}

  \caption{\texttt{block\_header\_t} 与用户数据区的内存布局}
  \label{fig:block-header-layout}
\end{figure}

\texttt{kmalloc} 返回 \texttt{block + HEADER\_SIZE}；\texttt{kfree} 时往回退 16 字节即找到 header。
32 位下 $\texttt{sizeof} = 4{+}4{+}4{+}4 = 16$，天然 8 字节对齐。

% ---------------------------------------------------------------------------
\subsection{分配：遍历 + 拆分}

\texttt{heap\_first\_fit\_alloc(size)} 步骤：
① 对齐 \texttt{ALIGN8(size)}；
② 从头遍历找第一个 \texttt{is\_free \&\& size >= need}；
③ 若块比请求大得多（余 $>$ \texttt{HEADER\_SIZE + 8}）则拆分；
④ \texttt{is\_free = 0}，返回数据区指针。

\begin{figure}[H]
  \centering
  % figures/ch10/fig04-block-split.tex — auto-extracted from chapter source
\begin{tikzpicture}[cell/.style={draw, minimum height=0.65cm, font=\ttfamily\small, anchor=west}]
  \node[font=\bfseries\small, anchor=west] at (-0.5, 1.2) {拆分前：};
  \node[cell, minimum width=1.8cm, fill=red!12]  at (0, 0.5) {\small hdr};
  \node[cell, minimum width=7cm, fill=green!12]   at (1.8, 0.5) {\small 大空闲块};
  \node[font=\bfseries\small, anchor=west] at (-0.5, -0.3) {拆分后：};
  \node[cell, minimum width=1.8cm, fill=red!12]   at (0, -1) {\small hdr};
  \node[cell, minimum width=2.5cm, fill=orange!15] at (1.8, -1) {\small 已分配};
  \node[cell, minimum width=1.8cm, fill=red!12]   at (4.3, -1) {\small new};
  \node[cell, minimum width=2.7cm, fill=green!12]  at (6.1, -1) {\small 新空闲};
  \draw[{|}-{|}, thick] (1.8, -1.9) -- (4.3, -1.9);
  \node[font=\scriptsize, below] at (3.05, -1.9) {\texttt{size}};
  \draw[{|}-{|}, thick] (6.1, -1.9) -- (8.8, -1.9);
  \node[font=\scriptsize, below] at (7.45, -1.9) {remaining};
\end{tikzpicture}

  \caption{First-Fit 分配时的块拆分}
  \label{fig:block-split}
\end{figure}

核心拆分代码：

\codefile{src/kernel/mm/heap\_first\_fit.c}
\begin{lstlisting}[style=cstyle]
if (curr->size - size > HEADER_SIZE + 8) {
    block_header_t* new_block =
        (block_header_t*)((uint8_t*)curr + HEADER_SIZE + size);
    new_block->size    = curr->size - size - HEADER_SIZE;
    new_block->is_free = 1;
    new_block->next = curr->next;  new_block->prev = curr;
    if (curr->next) curr->next->prev = new_block;
    curr->size = size;
    curr->next = new_block;
}
curr->is_free = 0;
return (void*)((uint8_t*)curr + HEADER_SIZE);
\end{lstlisting}

\begin{thinkbox}
拆分阈值 \texttt{HEADER\_SIZE + 8} 的含义：剩余空间必须够放 header（16\,B）加最小数据（8\,B），
否则拆出的碎片永远不可用。若改成 \texttt{HEADER\_SIZE + 1}，返回的指针将不满足 8 字节对齐——
许多 CPU 指令（如 SSE）会因此触发异常。
\end{thinkbox}

% ---------------------------------------------------------------------------
\subsection{释放与合并}

\texttt{heap\_first\_fit\_free(ptr)} 步骤：
① \texttt{block = ptr - HEADER\_SIZE}；
② 范围检查；
③ \texttt{is\_free = 1}；
④ 向后合并（while \texttt{next} 也 free 就 \texttt{merge\_with\_next}）；
⑤ 向前合并（若 \texttt{prev} 也 free 让其吞并自己）。

\begin{figure}[H]
  \centering
  % figures/ch10/fig05-block-coalesce.tex — auto-extracted from chapter source
\begin{tikzpicture}[cell/.style={draw, minimum height=0.65cm, font=\ttfamily\small, anchor=west}]
  \node[font=\bfseries\small, anchor=west] at (-0.5, 1.2) {合并前：};
  \node[cell, minimum width=1.4cm, fill=red!12]   at (0, 0.5) {\small hdr};
  \node[cell, minimum width=2.2cm, fill=green!12]  at (1.4, 0.5) {\small FREE};
  \node[cell, minimum width=1.4cm, fill=red!12]   at (3.6, 0.5) {\small hdr};
  \node[cell, minimum width=2.2cm, fill=green!12]  at (5.0, 0.5) {\small FREE};
  \node[cell, minimum width=1.4cm, fill=red!12]   at (7.2, 0.5) {\small hdr};
  \node[cell, minimum width=1.4cm, fill=orange!15] at (8.6, 0.5) {\small USED};
  \node[font=\bfseries\small, anchor=west] at (-0.5, -0.3) {合并后：};
  \node[cell, minimum width=1.4cm, fill=red!12]   at (0, -1) {\small hdr};
  \node[cell, minimum width=5.8cm, fill=green!12]  at (1.4, -1) {\small 合并后大块};
  \node[cell, minimum width=1.4cm, fill=red!12]   at (7.2, -1) {\small hdr};
  \node[cell, minimum width=1.4cm, fill=orange!15] at (8.6, -1) {\small USED};
  \draw[-{Stealth[length=3pt]}, dashed, gray] (4.0, 0.15) -- (4.0, -0.65);
  \node[font=\scriptsize\color{gray}] at (4.0, -0.25) {被吞并};
\end{tikzpicture}

  \caption{相邻空闲块的合并}
  \label{fig:block-coalesce}
\end{figure}

\codefile{src/kernel/mm/heap\_first\_fit.c}
\begin{lstlisting}[style=cstyle]
static void merge_with_next(block_header_t* block) {
    block_header_t* victim = block->next;
    block->size += HEADER_SIZE + victim->size;
    block->next  = victim->next;
    if (victim->next) victim->next->prev = block;
}
\end{lstlisting}

\begin{whybox}
为什么需要双向链表？

单向链表向后合并（吃 next）没问题，但向前合并需要从 block 找前驱——
单向链表要 $O(n)$ 遍历。双向链表多花 4\,B 存 \texttt{prev}，换来
$O(1)$ 向前合并，是经典的空间换时间。
\end{whybox}

% ============================================================================

